\chapter{Perzeptrons}

\section{Ziel dieses Kapitels}

In den vorherigen Kapiteln haben wir lineare Regression (linear regression) und logistische Regression (logistic regression) betrachtet.
Beide Modelle verwenden eine lineare Kombination der Eingaben:

\[
    a
    =
    \vec{w}^{\top}\vec{x}.
\]

Bei linearer Regression wird daraus direkt ein kontinuierlicher Wert vorhergesagt.
Bei logistischer Regression wird der lineare Score durch die Sigmoid-Funktion (sigmoid function) in eine Wahrscheinlichkeit umgewandelt.

In diesem Kapitel betrachten wir das Perzeptron (perceptron).
Das Perzeptron ist ein frühes und sehr wichtiges Modell für binäre Klassifikation.
Es ist außerdem ein historischer Vorläufer neuronaler Netze (neural networks).

\begin{conceptbox}
Ein Perzeptron (perceptron) ist ein einfacher linearer Klassifikator.

Es berechnet zuerst einen linearen Score

\[
    a
    =
    \vec{w}^{\top}\vec{x}
\]

und trifft danach mit einer Schwellenfunktion (threshold function) eine binäre Entscheidung.
\end{conceptbox}

Nach diesem Kapitel solltest du folgende Fragen beantworten können:

\begin{itemize}
    \item Was ist ein Perzeptron (perceptron)?
    \item Wie hängt das Perzeptron mit einem künstlichen Neuron zusammen?
    \item Was ist eine Schwellenfunktion (threshold function)?
    \item Was ist eine Aktivierungsfunktion (activation function)?
    \item Was ist eine Entscheidungsgrenze (decision boundary)?
    \item Was bedeutet lineare Separierbarkeit (linear separability)?
    \item Wie funktioniert die Perzeptron-Lernregel (perceptron learning rule)?
    \item Warum konvergiert das Perzeptron nur bei linear separierbaren Daten?
    \item Warum kann ein einzelnes Perzeptron das XOR-Problem nicht lösen?
    \item Wie hängen Perzeptron, logistische Regression und neuronale Netze zusammen?
\end{itemize}

\section{Motivation}

Wir betrachten ein binäres Klassifikationsproblem.
Die Zielwerte seien

\[
    t_n \in \{-1,+1\}.
\]

Die Eingaben seien

\[
    \vec{x}_n \in \mathbb{R}^D.
\]

Das Ziel ist, eine Funktion zu lernen, die für eine neue Eingabe \(\vec{x}\) eine Klasse vorhersagt:

\[
    y(\vec{x})
    \in
    \{-1,+1\}.
\]

Ein sehr einfacher Ansatz ist:
Wir berechnen zuerst einen linearen Score

\[
    a
    =
    \vec{w}^{\top}\vec{x}
\]

und entscheiden dann anhand des Vorzeichens von \(a\).

Wenn

\[
    a \geq 0,
\]

sagen wir Klasse \(+1\) voraus.
Wenn

\[
    a < 0,
\]

sagen wir Klasse \(-1\) voraus.

\begin{definitionbox}
Ein binärer linearer Klassifikator (binary linear classifier) entscheidet anhand des Vorzeichens eines linearen Scores:

\[
    y(\vec{x})
    =
    \operatorname{sign}(\vec{w}^{\top}\vec{x}).
\]
\end{definitionbox}

\section{Das Perzeptron-Modell}

Das Perzeptron besteht aus zwei Schritten:

\begin{enumerate}
    \item Lineare Kombination der Eingabe:
    \[
        a
        =
        \vec{w}^{\top}\vec{x}.
    \]

    \item Anwendung einer Schwellenfunktion:
    \[
        y
        =
        h(a).
    \]
\end{enumerate}

Für Klassen

\[
    y \in \{-1,+1\}
\]

verwendet man oft die Vorzeichenfunktion:

\[
    h(a)
    =
    \operatorname{sign}(a).
\]

Damit ist:

\[
    y(\vec{x},\vec{w})
    =
    \operatorname{sign}(\vec{w}^{\top}\vec{x}).
\]

\begin{definitionbox}
Das Perzeptron (perceptron) ist definiert durch

\[
    y(\vec{x},\vec{w})
    =
    \operatorname{sign}(\vec{w}^{\top}\vec{x}).
\]

Dabei ist \(\vec{x}\) der Eingabevektor, \(\vec{w}\) der Gewichtsvektor und \(y\) die vorhergesagte Klasse.
\end{definitionbox}

Die Vorzeichenfunktion ist:

\[
    \operatorname{sign}(a)
    =
    \begin{cases}
        +1, & a \geq 0, \\
        -1, & a < 0.
    \end{cases}
\]

\begin{notebox}
Manche Darstellungen verwenden Klassen

\[
    t \in \{0,1\}.
\]

Dann verwendet man eine Schrittfunktion, die Werte \(0\) und \(1\) ausgibt.
In diesem Kapitel verwenden wir meistens

\[
    t \in \{-1,+1\},
\]

weil die Perzeptron-Lernregel damit besonders einfach aussieht.
\end{notebox}

\section{Bias-Term}

Wie bei linearer und logistischer Regression ist ein Bias-Term wichtig.
Ohne Bias hätte das Modell die Form:

\[
    a
    =
    \sum_{j=1}^{D}
    w_jx_j.
\]

Die Entscheidungsgrenze müsste dann durch den Ursprung gehen.

Mit Bias-Term:

\[
    a
    =
    w_0
    +
    \sum_{j=1}^{D}
    w_jx_j.
\]

Wir können den Bias wieder in den Eingabevektor einbauen:

\[
    \tilde{\vec{x}}
    =
    \begin{pmatrix}
        1 \\
        x_1 \\
        \vdots \\
        x_D
    \end{pmatrix},
    \qquad
    \vec{w}
    =
    \begin{pmatrix}
        w_0 \\
        w_1 \\
        \vdots \\
        w_D
    \end{pmatrix}.
\]

Dann gilt:

\[
    a
    =
    \vec{w}^{\top}\tilde{\vec{x}}.
\]

\begin{conceptbox}
Der Bias-Term \(w_0\) verschiebt die Entscheidungsgrenze.

Ohne Bias muss die Entscheidungsgrenze durch den Ursprung gehen.
Mit Bias kann sie frei verschoben werden.
\end{conceptbox}

Im Folgenden schreiben wir oft einfach \(\vec{x}\), obwohl der Bias-Eintrag \(1\) bereits enthalten sein kann.

\section{Interpretation als künstliches Neuron}

Das Perzeptron kann als einfaches künstliches Neuron (artificial neuron) interpretiert werden.

Ein Neuron erhält Eingaben

\[
    x_1,x_2,\dots,x_D.
\]

Jede Eingabe wird mit einem Gewicht multipliziert:

\[
    w_1x_1,w_2x_2,\dots,w_Dx_D.
\]

Dann werden die gewichteten Eingaben aufsummiert:

\[
    a
    =
    w_0
    +
    \sum_{j=1}^{D}
    w_jx_j.
\]

Danach wird eine Aktivierungsfunktion (activation function) angewendet:

\[
    y
    =
    h(a).
\]

\begin{definitionbox}
Ein künstliches Neuron (artificial neuron) berechnet typischerweise

\[
    a
    =
    \sum_{j=1}^{D}
    w_jx_j
    +
    w_0
\]

und danach

\[
    y
    =
    h(a).
\]

Dabei ist \(h\) eine Aktivierungsfunktion (activation function).
\end{definitionbox}

Beim Perzeptron ist die Aktivierungsfunktion eine harte Schwellenfunktion:

\[
    h(a)
    =
    \operatorname{sign}(a).
\]

\section{Entscheidungsgrenze}

Das Perzeptron entscheidet nach dem Vorzeichen von

\[
    a
    =
    \vec{w}^{\top}\vec{x}.
\]

Die Entscheidungsgrenze (decision boundary) liegt dort, wo das Modell genau zwischen den Klassen steht:

\[
    \vec{w}^{\top}\vec{x}=0.
\]

\begin{definitionbox}
Die Entscheidungsgrenze eines Perzeptrons ist die Menge aller Punkte \(\vec{x}\), für die gilt:

\[
    \vec{w}^{\top}\vec{x}=0.
\]
\end{definitionbox}

Für zwei Eingabedimensionen mit Bias gilt:

\[
    \vec{x}
    =
    \begin{pmatrix}
        1 \\
        x_1 \\
        x_2
    \end{pmatrix}.
\]

Dann ist:

\[
    \vec{w}^{\top}\vec{x}
    =
    w_0+w_1x_1+w_2x_2.
\]

Die Entscheidungsgrenze ist:

\[
    w_0+w_1x_1+w_2x_2=0.
\]

Falls

\[
    w_2 \neq 0,
\]

kann man nach \(x_2\) auflösen:

\[
    x_2
    =
    -\frac{w_0}{w_2}
    -
    \frac{w_1}{w_2}x_1.
\]

Das ist eine Gerade.

\begin{conceptbox}
Ein einzelnes Perzeptron ist ein linearer Klassifikator.

In zwei Dimensionen ist die Entscheidungsgrenze eine Gerade.
In drei Dimensionen ist sie eine Ebene.
In \(D\) Dimensionen ist sie eine Hyperebene.
\end{conceptbox}

\section{Geometrische Interpretation des Gewichtsvektors}

Die Entscheidungsgrenze ist:

\[
    \vec{w}^{\top}\vec{x}=0.
\]

Der Gewichtsvektor \(\vec{w}\) steht senkrecht auf dieser Entscheidungsgrenze.
Er zeigt in die Richtung, in der der Score

\[
    a
    =
    \vec{w}^{\top}\vec{x}
\]

am schnellsten größer wird.

Wenn

\[
    \vec{w}^{\top}\vec{x} > 0,
\]

liegt \(\vec{x}\) auf der positiven Seite der Grenze.
Wenn

\[
    \vec{w}^{\top}\vec{x} < 0,
\]

liegt \(\vec{x}\) auf der negativen Seite.

\begin{notebox}
Die Länge von \(\vec{w}\) beeinflusst beim reinen Perzeptron nicht die vorhergesagte Klasse, solange das Vorzeichen gleich bleibt.

Wenn man \(\vec{w}\) mit einer positiven Konstanten \(c>0\) multipliziert, gilt:

\[
    \operatorname{sign}(c\vec{w}^{\top}\vec{x})
    =
    \operatorname{sign}(\vec{w}^{\top}\vec{x}).
\]
\end{notebox}

\section{Lineare Separierbarkeit}

Ein Datensatz ist linear separierbar, wenn es eine Hyperebene gibt, die die beiden Klassen perfekt trennt.

\begin{definitionbox}
Ein Datensatz

\[
    \mathcal{D}
    =
    \{(\vec{x}_1,t_1),\dots,(\vec{x}_N,t_N)\}
\]

mit

\[
    t_n \in \{-1,+1\}
\]

heißt linear separierbar (linearly separable), wenn es einen Gewichtsvektor \(\vec{w}\) gibt, sodass für alle \(n\) gilt:

\[
    t_n(\vec{w}^{\top}\vec{x}_n) > 0.
\]
\end{definitionbox}

Warum ist diese Bedingung sinnvoll?

Wenn

\[
    t_n=+1,
\]

dann brauchen wir

\[
    \vec{w}^{\top}\vec{x}_n > 0.
\]

Dann ist:

\[
    t_n(\vec{w}^{\top}\vec{x}_n)
    =
    +1\cdot(\vec{w}^{\top}\vec{x}_n)
    >
    0.
\]

Wenn

\[
    t_n=-1,
\]

dann brauchen wir

\[
    \vec{w}^{\top}\vec{x}_n < 0.
\]

Dann ist:

\[
    t_n(\vec{w}^{\top}\vec{x}_n)
    =
    -1\cdot(\vec{w}^{\top}\vec{x}_n)
    >
    0.
\]

\begin{conceptbox}
Die Bedingung

\[
    t_n(\vec{w}^{\top}\vec{x}_n)>0
\]

bedeutet:
Der Datenpunkt \(\vec{x}_n\) liegt auf der richtigen Seite der Entscheidungsgrenze.
\end{conceptbox}

\section{Korrekt und falsch klassifizierte Punkte}

Die Perzeptron-Vorhersage ist:

\[
    y_n
    =
    \operatorname{sign}(\vec{w}^{\top}\vec{x}_n).
\]

Ein Punkt ist korrekt klassifiziert, wenn

\[
    y_n=t_n.
\]

Für \(t_n \in \{-1,+1\}\) ist das äquivalent zu:

\[
    t_n(\vec{w}^{\top}\vec{x}_n)>0.
\]

Ein Punkt ist falsch klassifiziert, wenn

\[
    y_n\neq t_n.
\]

Das ist äquivalent zu:

\[
    t_n(\vec{w}^{\top}\vec{x}_n)\leq 0.
\]

\begin{definitionbox}
Die Menge der falsch klassifizierten Punkte ist

\[
    \mathcal{M}
    =
    \left\{
        n
        \mid
        t_n(\vec{w}^{\top}\vec{x}_n)\leq 0
    \right\}.
\]
\end{definitionbox}

\section{Perzeptron-Kriterium}

Um ein Perzeptron zu trainieren, braucht man eine Fehlerfunktion.
Eine Möglichkeit ist das Perzeptron-Kriterium (perceptron criterion).

Für falsch klassifizierte Punkte ist

\[
    t_n(\vec{w}^{\top}\vec{x}_n)\leq 0.
\]

Wir definieren den Fehler:

\[
    E_P(\vec{w})
    =
    -\sum_{n\in\mathcal{M}}
    t_n
    \vec{w}^{\top}\vec{x}_n.
\]

\begin{definitionbox}
Das Perzeptron-Kriterium (perceptron criterion) ist

\[
    E_P(\vec{w})
    =
    -\sum_{n\in\mathcal{M}}
    t_n
    \vec{w}^{\top}\vec{x}_n,
\]

wobei \(\mathcal{M}\) die Menge der falsch klassifizierten Punkte ist.
\end{definitionbox}

Für einen falsch klassifizierten Punkt gilt:

\[
    t_n\vec{w}^{\top}\vec{x}_n\leq 0.
\]

Daher ist

\[
    -t_n\vec{w}^{\top}\vec{x}_n\geq 0.
\]

Das Perzeptron-Kriterium ist also nichtnegativ:

\[
    E_P(\vec{w})\geq 0.
\]

Wenn alle Punkte korrekt klassifiziert sind, ist

\[
    \mathcal{M}=\emptyset
\]

und damit

\[
    E_P(\vec{w})=0.
\]

\begin{conceptbox}
Das Perzeptron-Kriterium bestraft nur falsch klassifizierte Punkte.

Korrekt klassifizierte Punkte tragen nicht zum Fehler bei.
\end{conceptbox}

\section{Gradient des Perzeptron-Kriteriums}

Für einen falsch klassifizierten Datenpunkt ist der Beitrag zum Fehler:

\[
    E_n(\vec{w})
    =
    -t_n\vec{w}^{\top}\vec{x}_n.
\]

Der Gradient nach \(\vec{w}\) ist:

\[
    \nabla_{\vec{w}}E_n(\vec{w})
    =
    -t_n\vec{x}_n.
\]

Für alle falsch klassifizierten Punkte:

\[
    \nabla_{\vec{w}}E_P(\vec{w})
    =
    -\sum_{n\in\mathcal{M}}
    t_n\vec{x}_n.
\]

\begin{definitionbox}
Der Gradient des Perzeptron-Kriteriums ist

\[
    \nabla_{\vec{w}}E_P(\vec{w})
    =
    -\sum_{n\in\mathcal{M}}
    t_n\vec{x}_n.
\]
\end{definitionbox}

\section{Perzeptron-Lernregel}

Gradientenabstieg verwendet:

\[
    \vec{w}^{(\tau+1)}
    =
    \vec{w}^{(\tau)}
    -
    \eta
    \nabla_{\vec{w}}E(\vec{w}^{(\tau)}).
\]

Für einen einzelnen falsch klassifizierten Punkt gilt:

\[
    \nabla_{\vec{w}}E_n(\vec{w})
    =
    -t_n\vec{x}_n.
\]

Also:

\[
    \vec{w}^{(\tau+1)}
    =
    \vec{w}^{(\tau)}
    -
    \eta(-t_n\vec{x}_n).
\]

Damit:

\[
    \vec{w}^{(\tau+1)}
    =
    \vec{w}^{(\tau)}
    +
    \eta t_n\vec{x}_n.
\]

\begin{definitionbox}
Die Perzeptron-Lernregel (perceptron learning rule) lautet für einen falsch klassifizierten Punkt:

\[
    \vec{w}^{(\tau+1)}
    =
    \vec{w}^{(\tau)}
    +
    \eta t_n\vec{x}_n.
\]

Für korrekt klassifizierte Punkte wird \(\vec{w}\) nicht geändert.
\end{definitionbox}

Dabei ist

\[
    \eta>0
\]

die Lernrate (learning rate).

\section{Intuition der Lernregel}

Die Lernregel ist sehr intuitiv.

\subsection{Fall \(t_n=+1\)}

Wenn der wahre Zielwert

\[
    t_n=+1
\]

ist, aber der Punkt falsch klassifiziert wird, dann ist

\[
    \vec{w}^{\top}\vec{x}_n\leq 0.
\]

Das Modell gibt also zu kleinen oder negativen Score aus.
Die Lernregel ist:

\[
    \vec{w}^{(\tau+1)}
    =
    \vec{w}^{(\tau)}
    +
    \eta \vec{x}_n.
\]

Dadurch wird der Score für diesen Punkt größer, denn:

\[
    (\vec{w}^{(\tau)}+\eta\vec{x}_n)^{\top}\vec{x}_n
    =
    (\vec{w}^{(\tau)})^{\top}\vec{x}_n
    +
    \eta\vec{x}_n^{\top}\vec{x}_n.
\]

Da

\[
    \vec{x}_n^{\top}\vec{x}_n
    =
    \|\vec{x}_n\|^2
    \geq 0,
\]

wird der Score erhöht.

\subsection{Fall \(t_n=-1\)}

Wenn der wahre Zielwert

\[
    t_n=-1
\]

ist, aber der Punkt falsch klassifiziert wird, dann ist

\[
    \vec{w}^{\top}\vec{x}_n\geq 0.
\]

Das Modell gibt also einen zu großen Score aus.
Die Lernregel ist:

\[
    \vec{w}^{(\tau+1)}
    =
    \vec{w}^{(\tau)}
    -
    \eta \vec{x}_n.
\]

Dadurch wird der Score für diesen Punkt kleiner.

\begin{conceptbox}
Die Perzeptron-Lernregel verschiebt \(\vec{w}\) so, dass der falsch klassifizierte Punkt beim nächsten Mal eher auf der richtigen Seite der Entscheidungsgrenze liegt.
\end{conceptbox}

\section{Algorithmus des Perzeptrons}

Der Trainingsalgorithmus läuft iterativ über die Trainingsdaten.

\begin{definitionbox}
Perzeptron-Algorithmus:

\begin{enumerate}
    \item Initialisiere \(\vec{w}\), zum Beispiel mit \(\vec{w}=\vec{0}\).
    \item Wähle einen Datenpunkt \((\vec{x}_n,t_n)\).
    \item Berechne den Score:
    \[
        a_n=\vec{w}^{\top}\vec{x}_n.
    \]
    \item Prüfe, ob der Punkt falsch klassifiziert ist:
    \[
        t_na_n\leq 0.
    \]
    \item Falls ja, aktualisiere:
    \[
        \vec{w}
        \leftarrow
        \vec{w}
        +
        \eta t_n\vec{x}_n.
    \]
    \item Wiederhole, bis keine Fehler mehr auftreten oder eine maximale Anzahl an Iterationen erreicht ist.
\end{enumerate}
\end{definitionbox}

\begin{notebox}
Die Reihenfolge der Datenpunkte kann das Training beeinflussen.

In praktischen Varianten werden die Datenpunkte oft in zufälliger Reihenfolge durchlaufen.
\end{notebox}

\section{Perzeptron-Konvergenz}

Ein klassisches Resultat ist der Perzeptron-Konvergenzsatz.

\begin{conceptbox}
Wenn die Daten linear separierbar sind, dann findet der Perzeptron-Algorithmus nach endlich vielen Updates eine trennende Hyperebene.

Wenn die Daten nicht linear separierbar sind, muss der Algorithmus nicht konvergieren.
\end{conceptbox}

Der genaue Konvergenzbeweis benötigt Annahmen über den Abstand der Daten zur trennenden Hyperebene.
Die wichtige Intuition ist:

\[
    \text{linear separierbar}
    \quad
    \Rightarrow
    \quad
    \text{Perzeptron kann eine perfekte lineare Trennung finden}.
\]

Wenn die Daten aber nicht linear separierbar sind, gibt es keine Hyperebene, die alle Punkte korrekt trennt.
Dann kann das Perzeptron immer wieder Updates durchführen, ohne jemals alle Punkte korrekt zu klassifizieren.

\begin{examtipbox}
Wichtig für Prüfungen:

\begin{itemize}
    \item Das Perzeptron konvergiert bei linear separierbaren Daten.
    \item Das Perzeptron konvergiert nicht garantiert bei nicht separierbaren Daten.
    \item Ein einzelnes Perzeptron kann nur lineare Entscheidungsgrenzen darstellen.
\end{itemize}
\end{examtipbox}

\section{Beispiel für ein Update}

Wir betrachten einen Datenpunkt

\[
    \vec{x}
    =
    \begin{pmatrix}
        1 \\
        2
    \end{pmatrix},
\]

mit Zielwert

\[
    t=+1.
\]

Der aktuelle Gewichtsvektor sei

\[
    \vec{w}
    =
    \begin{pmatrix}
        -1 \\
        0
    \end{pmatrix}.
\]

Der Score ist:

\[
    a
    =
    \vec{w}^{\top}\vec{x}
    =
    (-1)\cdot 1 + 0\cdot 2
    =
    -1.
\]

Da

\[
    t a
    =
    (+1)(-1)
    =
    -1
    \leq 0,
\]

ist der Punkt falsch klassifiziert.

Mit Lernrate

\[
    \eta=1
\]

ist das Update:

\[
    \vec{w}_{\mathrm{neu}}
    =
    \vec{w}
    +
    \eta t\vec{x}.
\]

Einsetzen:

\[
    \vec{w}_{\mathrm{neu}}
    =
    \begin{pmatrix}
        -1 \\
        0
    \end{pmatrix}
    +
    1\cdot(+1)
    \begin{pmatrix}
        1 \\
        2
    \end{pmatrix}.
\]

Also:

\[
    \vec{w}_{\mathrm{neu}}
    =
    \begin{pmatrix}
        0 \\
        2
    \end{pmatrix}.
\]

\begin{examplebox}
Nach dem Update ist der neue Score:

\[
    a_{\mathrm{neu}}
    =
    \vec{w}_{\mathrm{neu}}^{\top}\vec{x}
    =
    0\cdot 1+2\cdot 2
    =
    4.
\]

Damit ist

\[
    t a_{\mathrm{neu}}
    =
    (+1)4
    =
    4
    >
    0.
\]

Der Punkt ist nun korrekt klassifiziert.
\end{examplebox}

\section{Vergleich mit logistischer Regression}

Perzeptron und logistische Regression sind eng verwandt.
Beide berechnen zuerst einen linearen Score:

\[
    a
    =
    \vec{w}^{\top}\vec{x}.
\]

Der Unterschied liegt in der Ausgabefunktion und im Trainingsziel.

Beim Perzeptron:

\[
    y
    =
    \operatorname{sign}(a).
\]

Bei logistischer Regression:

\[
    y
    =
    \sigma(a).
\]

\begin{definitionbox}
Perzeptron:

\[
    y
    =
    \operatorname{sign}(\vec{w}^{\top}\vec{x}).
\]

Logistische Regression:

\[
    p(t=1\mid\vec{x},\vec{w})
    =
    \sigma(\vec{w}^{\top}\vec{x}).
\]
\end{definitionbox}

\section{Harte und weiche Entscheidung}

Das Perzeptron trifft eine harte Entscheidung:

\[
    y \in \{-1,+1\}.
\]

Logistische Regression gibt eine weiche Wahrscheinlichkeit aus:

\[
    y \in (0,1).
\]

\begin{conceptbox}
Das Perzeptron sagt nur, welche Klasse vorhergesagt wird.

Logistische Regression sagt zusätzlich, wie sicher das Modell in Form einer Wahrscheinlichkeit ist.
\end{conceptbox}

Die Sigmoid-Funktion ist glatt und differenzierbar:

\[
    \sigma'(a)
    =
    \sigma(a)(1-\sigma(a)).
\]

Die Vorzeichenfunktion ist nicht glatt und an

\[
    a=0
\]

nicht differenzierbar.

Das macht logistische Regression für gradientenbasierte Optimierung angenehmer.

\section{Perzeptron-Verlust und Kreuzentropie}

Beim Perzeptron werden falsch klassifizierte Punkte bestraft:

\[
    E_P(\vec{w})
    =
    -\sum_{n\in\mathcal{M}}
    t_n\vec{w}^{\top}\vec{x}_n.
\]

Bei logistischer Regression wird die negative Log-Likelihood minimiert:

\[
    E_{\mathrm{CE}}(\vec{w})
    =
    -\sum_{n=1}^{N}
    \left[
        t_n\log y_n
        +
        (1-t_n)\log(1-y_n)
    \right].
\]

\begin{notebox}
Die beiden Modelle können ähnliche lineare Entscheidungsgrenzen lernen, aber die Optimierungsprobleme sind unterschiedlich.

Logistische Regression ist probabilistisch motiviert.
Das klassische Perzeptron ist eher algorithmisch und geometrisch motiviert.
\end{notebox}

\section{Mehrere Ausgabekonventionen}

Es gibt zwei häufige Kodierungen für binäre Klassen.

\subsection{Kodierung mit \(\{-1,+1\}\)}

Beim Perzeptron ist üblich:

\[
    t \in \{-1,+1\}.
\]

Dann ist die Vorhersage:

\[
    y
    =
    \operatorname{sign}(\vec{w}^{\top}\vec{x}).
\]

Die Korrektheitsbedingung ist:

\[
    t\vec{w}^{\top}\vec{x}>0.
\]

\subsection{Kodierung mit \(\{0,1\}\)}

Bei logistischer Regression ist üblich:

\[
    t \in \{0,1\}.
\]

Dann ist:

\[
    p(t=1\mid\vec{x})
    =
    \sigma(\vec{w}^{\top}\vec{x}).
\]

\begin{notebox}
Beide Kodierungen beschreiben binäre Klassifikation.
Man muss aber bei Formeln konsequent bleiben.

Für Perzeptron-Updates ist \(\{-1,+1\}\) oft einfacher.
Für Bernoulli-Likelihood und Kreuzentropie ist \(\{0,1\}\) oft einfacher.
\end{notebox}

\section{Limitierung: XOR-Problem}

Ein einzelnes Perzeptron kann nur linear separierbare Daten klassifizieren.
Das klassische Gegenbeispiel ist XOR.

Wir betrachten Eingaben

\[
    x_1,x_2 \in \{0,1\}.
\]

Die XOR-Funktion ist:

\[
\begin{array}{cc|c}
x_1 & x_2 & t \\
\hline
0 & 0 & 0 \\
0 & 1 & 1 \\
1 & 0 & 1 \\
1 & 1 & 0
\end{array}
\]

Die positiven Punkte sind:

\[
    (0,1)
    \quad
    \text{und}
    \quad
    (1,0).
\]

Die negativen Punkte sind:

\[
    (0,0)
    \quad
    \text{und}
    \quad
    (1,1).
\]

Diese Punkte können in der Ebene nicht durch eine einzige Gerade getrennt werden.

\begin{conceptbox}
Das XOR-Problem ist nicht linear separierbar.

Ein einzelnes Perzeptron kann XOR daher nicht lösen.
\end{conceptbox}

\section{Warum mehrschichtige Netze helfen}

Obwohl ein einzelnes Perzeptron nur lineare Entscheidungsgrenzen darstellen kann, können mehrere Perzeptrons kombiniert werden.

Ein mehrschichtiges neuronales Netz (multilayer neural network) verwendet versteckte Einheiten (hidden units).
Diese berechnen Zwischenrepräsentationen:

\[
    h_j
    =
    g(\vec{w}_j^{\top}\vec{x}),
\]

wobei \(g\) eine Aktivierungsfunktion ist.

Die Ausgabe kann dann aus diesen versteckten Einheiten berechnet werden:

\[
    y
    =
    f\left(
        \sum_j
        v_j h_j
    \right).
\]

\begin{conceptbox}
Mehrschichtige neuronale Netze können nichtlineare Entscheidungsgrenzen darstellen.

Die erste Schicht transformiert die Eingabe.
Die nächste Schicht kann im transformierten Raum linear trennen.
\end{conceptbox}

\section{Perzeptron als Baustein neuronaler Netze}

Ein künstliches neuronales Netz besteht aus vielen Einheiten, die ähnlich wie Perzeptrons aufgebaut sind:

\[
    a_j
    =
    \sum_i
    w_{ji}x_i
    +
    b_j,
\]

und

\[
    z_j
    =
    g(a_j).
\]

Dabei ist:

\begin{itemize}
    \item \(a_j\) die Aktivierung vor der Nichtlinearität,
    \item \(z_j\) die Ausgabe der Einheit,
    \item \(w_{ji}\) das Gewicht von Eingabe \(i\) zu Einheit \(j\),
    \item \(b_j\) der Bias,
    \item \(g\) die Aktivierungsfunktion.
\end{itemize}

\begin{definitionbox}
Eine neuronale Einheit (neural unit) berechnet

\[
    z_j
    =
    g
    \left(
        \sum_i
        w_{ji}x_i
        +
        b_j
    \right).
\]
\end{definitionbox}

Beim klassischen Perzeptron ist

\[
    g(a)=\operatorname{sign}(a).
\]

In modernen neuronalen Netzen verwendet man meistens glatte oder stückweise lineare Aktivierungsfunktionen, zum Beispiel Sigmoid, Tanh oder ReLU.

\section{Aktivierungsfunktionen}

Aktivierungsfunktionen (activation functions) machen neuronale Netze nichtlinear.
Ohne nichtlineare Aktivierungsfunktionen wäre die Verkettung mehrerer linearer Schichten wieder nur eine lineare Abbildung.

\subsection{Sigmoid}

Die Sigmoid-Funktion ist:

\[
    \sigma(a)
    =
    \frac{1}{1+\exp(-a)}.
\]

Sie gibt Werte in

\[
    (0,1)
\]

aus.

\subsection{Tanh}

Die Hyperbolic-Tangent-Funktion ist:

\[
    \tanh(a)
    =
    \frac{\exp(a)-\exp(-a)}
    {\exp(a)+\exp(-a)}.
\]

Sie gibt Werte in

\[
    (-1,1)
\]

aus.

\subsection{ReLU}

Die ReLU-Funktion (rectified linear unit) ist:

\[
    \operatorname{ReLU}(a)
    =
    \max(0,a).
\]

Also:

\[
    \operatorname{ReLU}(a)
    =
    \begin{cases}
        a, & a>0, \\
        0, & a\leq 0.
    \end{cases}
\]

\begin{conceptbox}
Nichtlineare Aktivierungsfunktionen sind entscheidend.

Ohne sie kann ein mehrschichtiges Netz keine komplexeren Funktionen darstellen als ein einzelnes lineares Modell.
\end{conceptbox}

\section{Warum lineare Schichten allein nicht reichen}

Angenommen, wir haben zwei lineare Schichten ohne Aktivierungsfunktion:

\[
    \vec{h}
    =
    W_1\vec{x},
\]

und

\[
    \vec{y}
    =
    W_2\vec{h}.
\]

Dann gilt:

\[
    \vec{y}
    =
    W_2W_1\vec{x}.
\]

Setze

\[
    W
    =
    W_2W_1.
\]

Dann:

\[
    \vec{y}
    =
    W\vec{x}.
\]

Das ist wieder nur eine lineare Abbildung.

\begin{examtipbox}
Mehrere lineare Schichten ohne Nichtlinearität sind äquivalent zu einer einzigen linearen Schicht.

Deshalb braucht ein neuronales Netz nichtlineare Aktivierungsfunktionen.
\end{examtipbox}

\section{Perzeptron und Feature-Transformation}

Wie bei logistischer Regression kann auch ein Perzeptron mit Basisfunktionen verwendet werden.
Statt direkt \(\vec{x}\) zu klassifizieren, verwendet man eine transformierte Darstellung:

\[
    \vec{\phi}(\vec{x}).
\]

Dann ist:

\[
    y
    =
    \operatorname{sign}
    \left(
        \vec{w}^{\top}\vec{\phi}(\vec{x})
    \right).
\]

Die Entscheidungsgrenze im transformierten Raum ist linear:

\[
    \vec{w}^{\top}\vec{\phi}(\vec{x})=0.
\]

Im ursprünglichen Eingaberaum kann sie nichtlinear sein.

\begin{examplebox}
Für eine zweidimensionale Eingabe

\[
    \vec{x}
    =
    \begin{pmatrix}
        x_1 \\
        x_2
    \end{pmatrix}
\]

kann man quadratische Features verwenden:

\[
    \vec{\phi}(\vec{x})
    =
    \begin{pmatrix}
        1 \\
        x_1 \\
        x_2 \\
        x_1^2 \\
        x_2^2 \\
        x_1x_2
    \end{pmatrix}.
\]

Dann kann das Perzeptron nichtlineare Grenzen wie Ellipsen oder Parabeln darstellen.
\end{examplebox}

\section{Online-Lernen}

Das klassische Perzeptron ist ein Online-Lernalgorithmus (online learning algorithm).
Das bedeutet:
Die Gewichte werden nach einzelnen Datenpunkten aktualisiert.

\[
    \vec{w}
    \leftarrow
    \vec{w}
    +
    \eta t_n\vec{x}_n.
\]

Diese Aktualisierung wird nur durchgeführt, wenn der aktuelle Datenpunkt falsch klassifiziert ist.

\begin{definitionbox}
Online-Lernen (online learning) bedeutet, dass das Modell schrittweise mit einzelnen Datenpunkten aktualisiert wird, statt immer den gesamten Datensatz gleichzeitig zu verwenden.
\end{definitionbox}

\begin{notebox}
Online-Lernen ist nützlich, wenn Daten nacheinander eintreffen oder wenn der Datensatz zu groß ist, um ihn vollständig auf einmal zu verarbeiten.
\end{notebox}

\section{Lernrate}

Die Lernrate (learning rate) \(\eta\) bestimmt die Schrittweite der Updates:

\[
    \vec{w}
    \leftarrow
    \vec{w}
    +
    \eta t_n\vec{x}_n.
\]

Wenn \(\eta\) groß ist, sind die Updates groß.
Wenn \(\eta\) klein ist, sind die Updates klein.

Beim klassischen Perzeptron ist die genaue Wahl von \(\eta\) weniger kritisch als bei vielen anderen Gradientenverfahren, weil nur das Vorzeichen von \(\vec{w}^{\top}\vec{x}\) zählt.
Für \(\eta>0\) ändert eine konstante Skalierung der Gewichte die Entscheidungsgrenze im Wesentlichen nicht.

\begin{notebox}
Trotzdem kann \(\eta\) den Trainingsverlauf beeinflussen, insbesondere bei Varianten des Perzeptrons oder bei nicht separierbaren Daten.
\end{notebox}

\section{Pocket-Algorithmus}

Wenn Daten nicht linear separierbar sind, konvergiert das klassische Perzeptron nicht garantiert.
Eine praktische Erweiterung ist der Pocket-Algorithmus (pocket algorithm).

Die Idee ist:
Man speichert während des Trainings den bisher besten Gewichtsvektor.

\begin{conceptbox}
Der Pocket-Algorithmus behält den besten bisher gefundenen Gewichtsvektor in der Tasche.

Auch wenn spätere Updates schlechter werden, kann man am Ende den besten gespeicherten Gewichtsvektor verwenden.
\end{conceptbox}

Formal:

\begin{enumerate}
    \item Trainiere wie beim Perzeptron.
    \item Nach bestimmten Schritten berechne die Anzahl der Fehlklassifikationen.
    \item Wenn der aktuelle Gewichtsvektor besser ist als der bisher beste, speichere ihn.
    \item Gib am Ende den besten gespeicherten Gewichtsvektor aus.
\end{enumerate}

\begin{notebox}
Der Pocket-Algorithmus löst das Problem nicht perfekt, aber er macht das Perzeptron bei nicht separierbaren Daten praktischer.
\end{notebox}

\section{Vergleich: Perzeptron, logistische Regression und lineare Regression}

\[
\begin{array}{c|c|c|c}
\text{Modell}
&
\text{Ausgabe}
&
\text{Zieltyp}
&
\text{Training}
\\
\hline
\text{Lineare Regression}
&
\vec{w}^{\top}\vec{x}
&
t\in\mathbb{R}
&
\text{Least Squares}
\\
\hline
\text{Logistische Regression}
&
\sigma(\vec{w}^{\top}\vec{x})
&
t\in\{0,1\}
&
\text{Kreuzentropie}
\\
\hline
\text{Perzeptron}
&
\operatorname{sign}(\vec{w}^{\top}\vec{x})
&
t\in\{-1,+1\}
&
\text{Perzeptron-Regel}
\end{array}
\]

\begin{conceptbox}
Alle drei Modelle verwenden einen linearen Score.

Der Unterschied liegt darin, wie der Score interpretiert wird und welche Fehlerfunktion optimiert wird.
\end{conceptbox}

\section{Mehrklassen-Perzeptron}

Das Perzeptron kann auch auf mehrere Klassen erweitert werden.
Sei

\[
    t \in \{1,2,\dots,K\}.
\]

Für jede Klasse \(k\) gibt es einen Gewichtsvektor \(\vec{w}_k\).
Der Score für Klasse \(k\) ist:

\[
    a_k
    =
    \vec{w}_k^{\top}\vec{x}.
\]

Die vorhergesagte Klasse ist:

\[
    \hat{t}
    =
    \arg\max_{k}
    a_k.
\]

\begin{definitionbox}
Beim Mehrklassen-Perzeptron (multiclass perceptron) gilt:

\[
    \hat{t}
    =
    \arg\max_{k}
    \vec{w}_k^{\top}\vec{x}.
\]
\end{definitionbox}

Wenn die wahre Klasse \(t\) ist und die vorhergesagte Klasse \(\hat{t}\) falsch ist, kann man aktualisieren:

\[
    \vec{w}_t
    \leftarrow
    \vec{w}_t+\eta\vec{x},
\]

und

\[
    \vec{w}_{\hat{t}}
    \leftarrow
    \vec{w}_{\hat{t}}-\eta\vec{x}.
\]

\begin{conceptbox}
Beim Mehrklassen-Perzeptron wird der Gewichtsvektor der richtigen Klasse in Richtung des Datenpunkts verschoben.

Der Gewichtsvektor der fälschlich vorhergesagten Klasse wird vom Datenpunkt weggeschoben.
\end{conceptbox}

\section{Perzeptron mit Margin}

Beim klassischen Perzeptron reicht es, wenn

\[
    t_n\vec{w}^{\top}\vec{x}_n>0.
\]

Man kann aber eine Sicherheitsmarge (margin) verlangen.
Zum Beispiel:

\[
    t_n\vec{w}^{\top}\vec{x}_n \geq \gamma,
\]

wobei

\[
    \gamma>0
\]

eine gewünschte Margin ist.

\begin{definitionbox}
Die Margin eines Datenpunkts bezüglich \(\vec{w}\) ist

\[
    t_n\vec{w}^{\top}\vec{x}_n.
\]

Eine positive Margin bedeutet korrekte Klassifikation.
Eine große positive Margin bedeutet, dass der Punkt mit Sicherheitsabstand korrekt klassifiziert wird.
\end{definitionbox}

Margin-Ideen führen später zu Modellen wie Support Vector Machines.

\begin{notebox}
Das klassische Perzeptron versucht nur, die Punkte korrekt zu klassifizieren.
Es maximiert nicht explizit die Margin.
\end{notebox}

\section{Zusammenhang zu Support Vector Machines}

Support Vector Machines (SVMs) sind ebenfalls lineare Klassifikatoren, haben aber ein anderes Trainingsziel.
Während das Perzeptron irgendeine trennende Hyperebene finden kann, sucht eine SVM eine Hyperebene mit möglichst großer Margin.

\[
    \text{Perzeptron}
    \quad
    \Rightarrow
    \quad
    \text{irgendeine trennende Hyperebene}.
\]

\[
    \text{SVM}
    \quad
    \Rightarrow
    \quad
    \text{trennende Hyperebene mit großer Margin}.
\]

\begin{conceptbox}
Das Perzeptron ist ein wichtiger Vorläufer vieler linearer Klassifikationsmethoden.

Margin-basierte Methoden wie Support Vector Machines verbessern die Idee, indem sie nicht nur korrekte Klassifikation, sondern auch robuste Trennung verlangen.
\end{conceptbox}

\section{Zusammenhang zu Backpropagation}

Das klassische Perzeptron verwendet eine nicht differenzierbare Schwellenfunktion:

\[
    \operatorname{sign}(a).
\]

Für moderne neuronale Netze braucht man aber meistens differenzierbare Aktivierungsfunktionen, damit Gradienten mit Backpropagation berechnet werden können.

Deshalb verwendet man in neuronalen Netzen Funktionen wie:

\[
    \sigma(a),
    \qquad
    \tanh(a),
    \qquad
    \operatorname{ReLU}(a).
\]

\begin{notebox}
ReLU ist bei \(a=0\) nicht differenzierbar, aber sie ist stückweise linear und funktioniert praktisch sehr gut.
In Implementierungen verwendet man an \(a=0\) einfach eine geeignete Subgradienten-Konvention.
\end{notebox}

\section{Praktische Aspekte}

\subsection{Feature-Skalierung}

Wie bei logistischer Regression kann Feature-Skalierung wichtig sein.
Wenn Merkmale sehr unterschiedliche Größenordnungen haben, dominieren Merkmale mit großen Zahlenwerten die Updates.

Eine typische Standardisierung ist:

\[
    x_j'
    =
    \frac{x_j-\mu_j}{\sigma_j}.
\]

Dabei werden \(\mu_j\) und \(\sigma_j\) auf den Trainingsdaten berechnet.

\subsection{Datenreihenfolge}

Da das Perzeptron online aktualisiert wird, kann die Reihenfolge der Datenpunkte den Trainingsverlauf beeinflussen.
Oft mischt man die Daten zufällig.

\subsection{Nicht separierbare Daten}

Bei nicht separierbaren Daten sollte man nicht erwarten, dass das klassische Perzeptron perfekt konvergiert.
Dann sind logistische Regression, regularisierte Modelle oder neuronale Netze oft stabiler.

\section{Häufige Fehler}

\begin{examtipbox}
Typische Fehler bei Perzeptrons:

\begin{enumerate}
    \item Perzeptron mit logistischer Regression verwechseln.
    \item Die Vorzeichenkodierung \(t\in\{-1,+1\}\) mit \(t\in\{0,1\}\) mischen.
    \item Den Bias-Term vergessen.
    \item Denken, ein einzelnes Perzeptron könne nichtlinear separierbare Probleme lösen.
    \item Die Entscheidungsgrenze mit der Aktivierungsfunktion verwechseln.
    \item Die Perzeptron-Lernregel mit dem Gradienten der Kreuzentropie verwechseln.
    \item Vergessen, dass das Perzeptron nur bei linear separierbaren Daten garantiert konvergiert.
    \item Das XOR-Problem als linear separierbar ansehen.
    \item Mehrere lineare Schichten ohne Aktivierungsfunktion für ein echtes tiefes Netz halten.
    \item Nichtlineare Aktivierungsfunktionen in neuronalen Netzen unterschätzen.
\end{enumerate}
\end{examtipbox}

\section{Mini-Aufgaben}

\subsection{Aufgabe 1: Perzeptron-Vorhersage}

Gegeben sei

\[
    \vec{w}
    =
    \begin{pmatrix}
        -1 \\
        2
    \end{pmatrix},
    \qquad
    \vec{x}
    =
    \begin{pmatrix}
        1 \\
        1
    \end{pmatrix}.
\]

Berechne die Perzeptron-Vorhersage

\[
    y
    =
    \operatorname{sign}(\vec{w}^{\top}\vec{x}).
\]

\begin{examplebox}
Zuerst:

\[
    \vec{w}^{\top}\vec{x}
    =
    (-1)\cdot 1 + 2\cdot 1
    =
    1.
\]

Da

\[
    1>0,
\]

gilt:

\[
    y=+1.
\]
\end{examplebox}

\subsection{Aufgabe 2: Falsch klassifiziert oder korrekt?}

Gegeben sei

\[
    t=-1
\]

und

\[
    \vec{w}^{\top}\vec{x}=2.
\]

Ist der Punkt korrekt klassifiziert?

\begin{examplebox}
Ein Punkt ist korrekt klassifiziert, wenn

\[
    t\vec{w}^{\top}\vec{x}>0.
\]

Hier:

\[
    t\vec{w}^{\top}\vec{x}
    =
    (-1)\cdot 2
    =
    -2.
\]

Da

\[
    -2 \leq 0,
\]

ist der Punkt falsch klassifiziert.
\end{examplebox}

\subsection{Aufgabe 3: Perzeptron-Update}

Gegeben sei

\[
    \vec{w}
    =
    \begin{pmatrix}
        0 \\
        1
    \end{pmatrix},
    \qquad
    \vec{x}
    =
    \begin{pmatrix}
        1 \\
        2
    \end{pmatrix},
    \qquad
    t=-1,
\]

und

\[
    \eta=1.
\]

Angenommen, der Punkt ist falsch klassifiziert.
Berechne das Update.

\begin{examplebox}
Die Lernregel ist:

\[
    \vec{w}_{\mathrm{neu}}
    =
    \vec{w}
    +
    \eta t\vec{x}.
\]

Einsetzen:

\[
    \vec{w}_{\mathrm{neu}}
    =
    \begin{pmatrix}
        0 \\
        1
    \end{pmatrix}
    +
    1\cdot(-1)
    \begin{pmatrix}
        1 \\
        2
    \end{pmatrix}.
\]

Also:

\[
    \vec{w}_{\mathrm{neu}}
    =
    \begin{pmatrix}
        0 \\
        1
    \end{pmatrix}
    -
    \begin{pmatrix}
        1 \\
        2
    \end{pmatrix}
    =
    \begin{pmatrix}
        -1 \\
        -1
    \end{pmatrix}.
\]
\end{examplebox}

\subsection{Aufgabe 4: Entscheidungsgrenze}

Gegeben sei ein Perzeptron mit

\[
    a
    =
    w_0+w_1x_1+w_2x_2.
\]

Die Gewichte seien

\[
    w_0=-2,
    \qquad
    w_1=1,
    \qquad
    w_2=1.
\]

Bestimme die Entscheidungsgrenze.

\begin{examplebox}
Die Entscheidungsgrenze ist:

\[
    a=0.
\]

Also:

\[
    -2+x_1+x_2=0.
\]

Damit:

\[
    x_2
    =
    2-x_1.
\]

Das ist eine Gerade.
\end{examplebox}

\subsection{Aufgabe 5: XOR}

Warum kann ein einzelnes Perzeptron XOR nicht lösen?

\begin{examplebox}
Ein einzelnes Perzeptron hat eine lineare Entscheidungsgrenze:

\[
    \vec{w}^{\top}\vec{x}=0.
\]

Für \(D=2\) ist das eine Gerade.

Die XOR-Datenpunkte sind aber nicht durch eine einzige Gerade trennbar.
Daher kann ein einzelnes Perzeptron XOR nicht lösen.
\end{examplebox}

\section{Zusammenfassung}

\begin{itemize}
    \item Ein Perzeptron (perceptron) ist ein linearer Klassifikator.
    \item Es berechnet einen Score:
    \[
        a
        =
        \vec{w}^{\top}\vec{x}.
    \]
    \item Die Ausgabe ist:
    \[
        y
        =
        \operatorname{sign}(a).
    \]
    \item Mit Klassen \(t\in\{-1,+1\}\) ist ein Punkt korrekt klassifiziert, wenn:
    \[
        t\vec{w}^{\top}\vec{x}>0.
    \]
    \item Die Entscheidungsgrenze ist:
    \[
        \vec{w}^{\top}\vec{x}=0.
    \]
    \item Ein einzelnes Perzeptron kann nur linear separierbare Daten perfekt trennen.
    \item Das Perzeptron-Kriterium ist:
    \[
        E_P(\vec{w})
        =
        -\sum_{n\in\mathcal{M}}
        t_n\vec{w}^{\top}\vec{x}_n.
    \]
    \item Die Perzeptron-Lernregel ist:
    \[
        \vec{w}
        \leftarrow
        \vec{w}
        +
        \eta t_n\vec{x}_n.
    \]
    \item Das Perzeptron konvergiert garantiert, wenn die Daten linear separierbar sind.
    \item Das XOR-Problem ist nicht linear separierbar und kann daher nicht von einem einzelnen Perzeptron gelöst werden.
    \item Mehrschichtige neuronale Netze verwenden Perzeptron-ähnliche Einheiten mit nichtlinearen Aktivierungsfunktionen.
    \item Nichtlineare Aktivierungsfunktionen sind notwendig, weil mehrere lineare Schichten ohne Nichtlinearität wieder nur eine lineare Abbildung ergeben.
\end{itemize}

\section{Typische Prüfungsfragen}

\begin{examtipbox}
Mögliche Prüfungsfragen zu diesem Kapitel:

\begin{enumerate}
    \item Definiere ein Perzeptron.
    \item Was ist der Unterschied zwischen Perzeptron und logistischer Regression?
    \item Welche Rolle spielt der Bias-Term?
    \item Was ist die Entscheidungsgrenze eines Perzeptrons?
    \item Warum ist ein einzelnes Perzeptron ein linearer Klassifikator?
    \item Was bedeutet lineare Separierbarkeit?
    \item Zeige, warum \(t_n\vec{w}^{\top}\vec{x}_n>0\) korrekte Klassifikation bedeutet.
    \item Definiere das Perzeptron-Kriterium.
    \item Leite den Gradienten des Perzeptron-Kriteriums für einen falsch klassifizierten Punkt her.
    \item Leite die Perzeptron-Lernregel her.
    \item Erkläre intuitiv, warum die Perzeptron-Lernregel sinnvoll ist.
    \item Wann konvergiert der Perzeptron-Algorithmus?
    \item Warum konvergiert er bei nicht separierbaren Daten nicht garantiert?
    \item Warum kann ein einzelnes Perzeptron XOR nicht lösen?
    \item Wie können Basisfunktionen einem Perzeptron nichtlineare Grenzen ermöglichen?
    \item Warum brauchen mehrschichtige neuronale Netze nichtlineare Aktivierungsfunktionen?
    \item Was ist der Unterschied zwischen Online-Lernen und Batch-Lernen?
    \item Was ist der Pocket-Algorithmus?
    \item Wie funktioniert das Mehrklassen-Perzeptron?
    \item Wie hängt das Perzeptron historisch und mathematisch mit neuronalen Netzen zusammen?
\end{enumerate}
\end{examtipbox}